\chapter{GIỚI THIỆU ĐỀ TÀI}
\section{Đặt vấn đề}

Trong những năm gần đây, hình thức thi trắc nghiệm đã trở thành phương thức đánh giá phổ biến tại các trường trung học, cao đẳng và đại học ở Việt Nam nhờ khả năng đánh giá nhanh, khách quan và phù hợp với các kỳ thi có số lượng thí sinh lớn. Tuy nhiên, khi số lượng bài thi ngày càng tăng, công tác chấm điểm và quản lý kết quả cũng đặt ra nhiều yêu cầu về tốc độ, độ chính xác và khả năng tự động hóa.

Hiện nay, các cơ sở giáo dục thường áp dụng một trong hai phương pháp chấm điểm chính.

Phương pháp thứ nhất là \textbf{chấm thủ công}. Giáo viên đối chiếu từng bài làm với đáp án và tính điểm bằng tay. Phương pháp này có ưu điểm là không yêu cầu thiết bị chuyên dụng nhưng tiêu tốn nhiều thời gian, đặc biệt đối với các lớp học hoặc kỳ thi có số lượng lớn bài làm. Quá trình chấm kéo dài cũng làm tăng khả năng xảy ra sai sót do nhầm lẫn hoặc mệt mỏi của người chấm, đồng thời gây khó khăn trong việc tổng hợp, thống kê và lưu trữ kết quả.

Phương pháp thứ hai là sử dụng \textbf{máy chấm phiếu trắc nghiệm OMR (Optical Mark Recognition)}. Đây là giải pháp được nhiều đơn vị tổ chức thi quy mô lớn lựa chọn nhờ tốc độ xử lý cao và độ chính xác tốt. Tuy nhiên, hệ thống OMR truyền thống yêu cầu sử dụng máy quét chuyên dụng có chi phí đầu tư lớn, thường dao động từ vài chục đến hàng trăm triệu đồng. Ngoài ra, các hệ thống này thường yêu cầu phiếu thi được in theo mẫu cố định, ảnh quét phải đạt chất lượng cao và việc vận hành phụ thuộc vào phòng máy hoặc thiết bị chuyên biệt, làm giảm tính linh hoạt trong các môi trường giảng dạy thông thường.

Trong bối cảnh điện thoại thông minh ngày càng được sử dụng rộng rãi và các kỹ thuật thị giác máy tính (Computer Vision) phát triển mạnh mẽ, việc tận dụng camera trên thiết bị di động để tự động nhận diện và chấm phiếu trắc nghiệm trở thành một hướng nghiên cứu có nhiều tiềm năng. Giải pháp này giúp giảm đáng kể chi phí triển khai, tăng tính cơ động và cho phép giáo viên thực hiện chấm điểm ở bất kỳ đâu mà không cần đầu tư hệ thống phần cứng chuyên dụng.

Tuy nhiên, việc xây dựng một hệ thống chấm điểm dựa trên ảnh chụp từ điện thoại cũng đặt ra nhiều thách thức kỹ thuật. Ảnh thu được trong thực tế thường xuất hiện hiện tượng nghiêng, méo phối cảnh, thay đổi điều kiện ánh sáng, bóng đổ hoặc chất lượng tô đáp án không đồng đều. Do đó, hệ thống cần có khả năng phát hiện chính xác vùng phiếu, hiệu chỉnh hình học, chuẩn hóa ảnh và nhận diện các ô tô với độ chính xác cao trong nhiều điều kiện khác nhau.

Xuất phát từ những yêu cầu trên, đề tài \textbf{"Xây dựng hệ thống web chấm điểm phiếu trắc nghiệm sử dụng kỹ thuật xử lý ảnh"} được thực hiện nhằm phát triển một giải pháp chấm điểm tự động hoạt động trên nền tảng web. Hệ thống cho phép giáo viên sử dụng camera của điện thoại hoặc máy tính để quét phiếu, tự động nhận diện mã số sinh viên, mã đề và đáp án, sau đó tính điểm và quản lý kết quả một cách nhanh chóng, chính xác mà không cần sử dụng máy chấm OMR chuyên dụng.

\section{Mục tiêu và phạm vi đề tài}
Đề tài hướng đến xây dựng một hệ thống web ứng dụng di động cho phép:

\begin{enumerate}
    \item \textbf{Tự động hóa hoàn toàn}: Giáo viên chỉ cần chụp ảnh phiếu bằng điện thoại, hệ thống xử lý và trả kết quả trong vòng vài giây.
    \item \textbf{Không phụ thuộc phần cứng chuyên dụng}: Chạy trên smartphone phổ thông, không cần máy quét.
    \item \textbf{Hỗ trợ nhiều loại phiếu}: Kiến trúc Form Profile linh hoạt cho phép cấu hình theo nhiều mẫu phiếu khác nhau.
    \item \textbf{Xử lý ảnh thực tế}: Chịu được ảnh chụp nghiêng, tô không rõ thông qua các thuật toán xử lý ảnh tiên tiến.
    \item \textbf{Quản lý theo bài thi}: Giáo viên tạo bài thi, quản lý nhiều mã đề, theo dõi lịch sử chấm và xuất báo cáo.
\end{enumerate}

Trong phạm vi của đề tài, hệ thống được thiết kế dưới dạng một ứng dụng web hoạt động ổn định trên môi trường mạng. Chức năng cốt lõi của hệ thống tập trung vào việc tự động chấm điểm các mẫu phiếu trắc nghiệm nhiều lựa chọn (hỗ trợ các dạng 4 đáp án A, B, C, D). Bên cạnh đó, ứng dụng được tích hợp khả năng nhận diện mã số sinh viên (MSSV) thông qua lưới ô tô số, đồng thời tự động trích xuất mã đề thi từ lưới ô tô số. Để tối ưu hóa quá trình nhập liệu, hệ thống cung cấp tính năng quét trực tiếp (live scanner) bằng camera hoạt động ngay trên trình duyệt web. Cuối cùng, toàn bộ kết quả chấm điểm có thể được trích xuất dễ dàng dưới các định dạng tiêu chuẩn như Excel và PDF, phục vụ hiệu quả cho công tác thống kê và báo cáo.
\section{Định hướng giải pháp}
Để đáp ứng các mục tiêu đã đặt ra và giải quyết triệt để những hạn chế của các phương pháp chấm điểm truyền thống, hệ thống được thiết kế theo kiến trúc ba tầng (Three-Tier Architecture) và phát triển dựa trên các công nghệ hiện đại, mang lại hiệu năng cao cùng khả năng mở rộng linh hoạt:

\begin{itemize}
    \item \textbf{Phần giao diện người dùng (Frontend):} Được phát triển bằng \textbf{React} kết hợp với \textbf{TypeScript}, cung cấp giao diện web động được tối ưu hóa mạnh mẽ cho thiết bị di động (mobile-first). Đặc biệt, tính năng Smart Camera Scanner sử dụng \textbf{WebRTC MediaDevices API}~\cite{webrtc} cho phép truy cập luồng video và phát hiện phiếu trắc nghiệm theo thời gian thực ngay trên trình duyệt mà không cần cài đặt ứng dụng native.

    \item \textbf{Phần xử lý logic (Backend):} Sử dụng \textbf{FastAPI}~\cite{fastapi} (Python) làm REST API server. Với cơ chế xử lý bất đồng bộ (Asynchronous) dựa trên nền tảng ASGI, FastAPI đảm bảo khả năng đáp ứng đồng thời nhiều yêu cầu xử lý ảnh với độ trễ thấp, đồng thời dễ dàng tích hợp với các thư viện đặc thù của Python.

    \item \textbf{Thị giác máy tính (Computer Vision):} Thư viện \textbf{OpenCV}~\cite{opencv,learning-opencv} đóng vai trò cốt lõi trong việc thực thi pipeline xử lý ảnh 15 bước. Các kỹ thuật như phát hiện cạnh Canny~\cite{canny}, biến đổi phối cảnh (Perspective Warp) và phân tích mật độ pixel được áp dụng để trích xuất và chuẩn hóa vùng dữ liệu phiếu từ các bức ảnh chụp trong điều kiện ánh sáng thực tế.

    \item \textbf{Cơ sở dữ liệu (Database):} Hệ thống sử dụng \textbf{PostgreSQL}~\cite{postgresql} làm cơ sở dữ liệu chính thông qua \textbf{SQLAlchemy ORM}~\cite{sqlalchemy}. PostgreSQL được lựa chọn nhờ khả năng hỗ trợ lưu trữ và truy vấn dữ liệu kiểu \texttt{JSON} một cách tự nhiên, cực kỳ phù hợp để lưu trữ cấu trúc linh hoạt của các bộ đáp án và lịch sử chấm điểm phức tạp của người dùng.
\end{itemize}
\section{Bố cục đồ án}
Báo cáo đồ án tốt nghiệp được tổ chức thành 7 chương với các nội dung chính như sau:

\begin{itemize}
    \item \textbf{Chương 1: Giới thiệu đề tài.} Trình bày lý do chọn đề tài, mục tiêu, phạm vi nghiên cứu và định hướng giải pháp công nghệ cho hệ thống.

    \item \textbf{Chương 2: Khảo sát và phân tích yêu cầu.} Trình bày kết quả khảo sát hiện trạng chấm điểm tại các cơ sở giáo dục và phân tích các hạn chế của những giải pháp hiện có. Xác định các yêu cầu của hệ thống thông qua biểu đồ use case, đặc tả chức năng chi tiết và các yêu cầu phi chức năng về hiệu năng, bảo mật, cũng như khả năng mở rộng.

    \item \textbf{Chương 3: Công nghệ sử dụng.} Trình bày các công nghệ được chọn, các lựa chọn thay thế tương đương và lý do lựa chọn theo yêu cầu kỹ thuật của hệ thống.

    \item \textbf{Chương 4: Xây dựng ứng dụng web OMR.} Trình bày thiết kế kiến trúc ba tầng và cấu trúc mã nguồn; thiết kế backend (sơ đồ gói, luồng phụ thuộc một chiều) và cơ sở dữ liệu (ERD, vai trò từng bảng); thiết kế giao diện frontend bao gồm biểu đồ máy trạng thái Smart Camera Scanner; các thành phần giao diện chi tiết, cấu hình Form Profile và ROI; quy trình triển khai và định hướng mở rộng lên môi trường đám mây; kết quả kiểm thử chức năng và hiệu năng.

    \item \textbf{Chương 5: Pipeline xử lý ảnh OMR và đánh giá.} Trình bày chi tiết pipeline xử lý ảnh 15 bước từ chuẩn hóa ảnh, xác định vùng đọc, giải mã MSSV và MCQ đến tính điểm và lưu kết quả; luồng chấm một phiếu và luồng chấm hàng loạt kèm biểu đồ trình tự; đánh giá định lượng hệ thống trên bộ dữ liệu thực tế ở nhiều điều kiện chụp khác nhau.

    \item \textbf{Chương 6: Các giải pháp và đóng góp nổi bật.} Đi sâu phân tích các giải pháp kỹ thuật mang tính đột phá của đồ án, bao gồm cơ chế MCQ Map Search Rescue, thuật toán tính điểm bong bóng tổng hợp và tính năng Smart Camera Scanner hoạt động trực tiếp trên trình duyệt.

    \item \textbf{Chương 7: Kết luận và hướng phát triển.} Tóm tắt những kết quả đã đạt được, đánh giá các điểm mạnh, điểm còn hạn chế của hệ thống và đề xuất các định hướng phát triển, mở rộng trong tương lai.
\end{itemize}
